\documentclass[../DoAn.tex]{subfiles}
\begin{document}

Chương 4 đã mô tả chi tiết quy trình lập trình, cài đặt thực tế các cấu phần mã nguồn từ lõi Spring Boot Backend, dịch vụ AI Flask, ứng dụng Flutter cho đến firmware nhúng ESP32. Chương 5 này tập trung trình bày các giải pháp kỹ thuật cốt lõi, mô hình kiến trúc an ninh và đóng góp công nghệ nổi bật mà đề tài đã nghiên cứu, hiện thực hóa thành công. Sự có mặt của chương này là cơ sở đặc biệt quan trọng để đánh giá năng lực tư duy, khả năng phân tích biện luận khoa học và tính sáng tạo của sinh viên trong suốt quá trình thực hiện đồ án tốt nghiệp cử nhân Công nghệ thông tin. Theo đúng quy chuẩn hướng dẫn của nhà trường, các đóng góp nổi bật trong chương sẽ được phân tách thành các mục nghiên cứu độc lập, phân rã chi tiết theo cấu trúc ba cấu phần bao gồm: (i) dẫn dắt, giới thiệu về bài toán và thách thức kỹ thuật đặt ra, (ii) giải pháp công nghệ cụ thể được thiết kế, triển khai, và (iii) kết quả định lượng đạt được trong môi trường thực nghiệm.

\section{Giải pháp kiến trúc xác thực hai lớp tự trị không kho ảnh tập trung}
\label{section:5.1}

\subsection{Dẫn dắt bài toán và thách thức kỹ thuật}

In các hệ thống kiểm soát ra vào và điểm danh sinh trắc học truyền thống, ban quản lý bắt buộc phải tổ chức một quy trình thu thập ảnh chân dung mẫu của toàn bộ nhân sự mới, sau đó tiến hành trích xuất và lưu trữ tập trung dữ liệu hình ảnh hoặc vector đặc trưng trong một cơ sở dữ liệu máy chủ trung tâm. Mô hình này đặt ra ba thách thức thực tế vô cùng nan giải:
\begin{enumerate}
    \item \textbf{Nguy cơ rủi ro pháp lý và bảo mật thông tin:} Việc lưu trữ tập trung kho ảnh sinh trắc học của hàng ngàn nhân sự tạo ra một điểm chết an ninh. Nếu máy chủ bị tấn công xâm nhập, toàn bộ kho ảnh gốc sẽ bị rò rỉ, vi phạm nghiêm trọng các quy định pháp lý về bảo vệ dữ liệu cá nhân.
    \item \textbf{Gánh nặng quy trình hành chính:} Tổ chức luôn tốn kém thời gian và nguồn lực để vận hành phân hệ đăng ký mẫu, chụp ảnh chuẩn hóa cho nhân sự mới mỗi khi có biến động nhân sự.
    \item \textbf{Sự sai lệch theo thời gian (Aging effect):} Ảnh mẫu thu thập thủ công từ nhiều năm trước có độ chính xác giảm dần do sự thay đổi tự nhiên trên khuôn mặt nhân sự, dẫn đến tỷ lệ từ chối sai (FRR) tăng cao.
\end{enumerate}

\subsection{Giải pháp kỹ thuật đề xuất}

Để giải quyết triệt để các thách thức trên, đề tài đề xuất giải pháp kiến trúc xác thực hai lớp tự trị phân tán (\textit{2-Layer Autonomous Authentication}), tận dụng trực tiếp hạ tầng định danh quốc gia sẵn có là thẻ Căn cước công dân gắn chip của Việt Nam để làm sạch kho dữ liệu. Cơ chế vận hành phân tách trách nhiệm được thiết kế tinh gọn như sau:
\begin{itemize}
    \item \textbf{Lớp xác thực 1 (Layer 1 - Phân tán on-device):} Vận hành cục bộ trực tiếp trên trạm đầu cuối di động Android thông qua ứng dụng Flutter mà không cần kết nối Internet. Ứng dụng thực hiện giao tiếp trường gần NFC để bóc tách ảnh chân dung gốc chất lượng cao do Bộ Công an chụp bảo chứng làm ảnh tham chiếu chuẩn. Đồng thời, thiết bị chụp ảnh selfie thực tế của người quét để chạy thuật toán trích xuất không gian hình học vector đặc trưng 512 chiều, so khớp qua khoảng cách Cosine similarity nhằm xác chuẩn chủ thẻ ngay tại chỗ theo ngưỡng an toàn diện hẹp 0.65.
    \item \textbf{Lớp xác thực 2 (Layer 2 - Tập trung chuyên sâu):} Khi lớp 1 đạt, ứng dụng Flutter đóng gói chuỗi dữ liệu mã hóa truyền trực tiếp trong phần thân yêu cầu gửi HTTP POST lên máy chủ Backend Spring Boot qua kênh mã hóa HTTPS bảo mật. Máy chủ trung tâm tiếp nhận, gọi không đồng bộ sang dịch vụ AI Python Flask thực hiện thuật toán đối sánh độ tương đồng toán học thông qua mô hình học sâu ArcFace. Đồng thời, Backend thực hiện đối chiếu thời gian thực với ma trận phân quyền truy cập \path{access_permissions} dựa trên mốc giờ kết thúc ca làm việc cá nhân (\texttt{end\_time}) nhằm tự động đóng phiên và làm sạch chuỗi trạng thái di chuyển tuần tự của nhân sự.
\end{itemize}

Bằng việc cấu hình bỏ qua bước đối soát cơ sở dữ liệu ảnh tập trung, hệ thống chuyển dịch hoàn toàn sang cơ chế kiểm tra chéo động: Xác thực thực thể thẻ vật lý kết hợp sinh trắc học của chính chủ sở hữu thẻ tại thời điểm quét.

\subsection{Kết quả thực nghiệm định lượng đạt được}

Giải pháp kiến trúc này loại bỏ hoàn toàn 100\% gánh nặng rủi ro pháp lý và chi phí xây dựng kho lưu trữ ảnh tập trung. Qua tổng số hơn 12,000 lượt thực nghiệm đối soát chéo hình ảnh chân dung giữa ảnh camera selfie thực tế và ảnh DG2 chất lượng cao bóc tách từ chip Căn cước công dân, hệ thống đo lường chính xác chỉ số tỷ lệ chấp nhận sai (\textit{False Acceptance Rate - FAR}) đạt ngưỡng cực thấp là 0.01\%, song song đó chỉ số tỷ lệ từ chối sai (\textit{False Rejection Rate - FRR}) đạt mức 1.15\% tại cấu hình điểm số ngưỡng khoảng cách hình học Cosine Similarity tối ưu đạt $0.65$. Mô hình mang lại điểm cân bằng an ninh lý tưởng: tối ưu hóa FPR xuống mức tối thiểu trong khi vẫn bảo đảm tốc độ tính toán ấn tượng, tạo hành lang an toàn dữ liệu tuyệt đối cho tổ chức.

\section{Giải pháp tối ưu hóa băng thông truyền thông ngoại vi qua Multi-Protocol Async Pipeline}
\label{section:5.2}

\subsection{Dẫn dắt bài toán và thách thức kỹ thuật}

Thách thức lớn thứ hai của đồ án là bài toán điều khiển phần cứng vật lý và đồng bộ giao diện Dashboard thời gian thực trong môi trường phân tán đa thiết bị. Nếu áp dụng các phương thức giao tiếp truyền thông truyền thống như HTTP Polling (bo mạch ESP32 liên tục gửi request xin lệnh sau mỗi 1 giây, hoặc trình duyệt React gửi request tải lại dữ liệu liên tục), hệ thống sẽ lập tức rơi vào trạng thái nghẽn mạch cục bộ khi mở rộng quy mô. Gánh nặng HTTP Request rác sẽ chiếm dụng toàn bộ tài nguyên bộ nhớ Thread Pool của máy chủ, gây phình trễ phản hồi vượt quá 10 giây và làm mất hoàn toàn tính chất giám sát trực tuyến của hệ thống kiểm soát ra vào an ninh.

\subsection{Giải pháp kỹ thuật đề xuất}

Đề tài giải quyết triệt để bài toán này bằng việc thiết kế và hiện thực hóa mô hình truyền thông bất đồng bộ đa giao thức \textit{Multi-Protocol Asynchronous Pipeline}, tích hợp song song các đường ống truyền dữ liệu chuyên dụng:
\begin{itemize}
    \item \textbf{Đường ống điều khiển nhúng ngoại vi (MQTT Channel):} Máy chủ Backend Spring Boot tích hợp thư viện nhúng không đồng bộ MQTT Client kết nối trực tiếp sang Broker HiveMQ Cloud thông qua cổng bảo mật SSL/TLS 8883. Ngay khi lõi Backend phê duyệt điểm danh thành công, một thông điệp nhỏ gọn \texttt{\{"action":"OPEN\_DOOR"\}} lập tức được xuất bản xuống chủ đề mạng \path{cccd/devices/command}. Mạch vi điều khiển ESP32 túc trực lắng nghe liên tục trên kết nối TCP kéo dài độc lập, tiếp nhận chuỗi JSON và kích hoạt chân GPIO điều khiển rơ-le mở khóa vật lý gần như tức thời.
    \item \textbf{Đường ống broadcast giao diện quản trị (WebSocket STOMP Channel):} Song song với luồng MQTT, Backend truyền thông điệp dạng gói tin STOMP đến các client trình duyệt React đang duy trì kết nối Socket dài hạn trên topic \path{/topic/attendance} vận hành tại cổng 3002. 
\end{itemize}

Để giải quyết bài toán mâu thuẫn giữa an ninh và lưu thông thực địa, đồ án triển khai Cơ chế chống quay vòng thẻ mềm tích hợp sinh trắc học dòng kép (Biometric-Driven Soft Mode). Khi phát hiện các lỗi nghịch đảo chuỗi di chuyển (như quẹt chiều RA liên tiếp hoặc thiếu dữ liệu VÀO), nếu trắc lượng khuôn mặt đạt chuẩn chính chủ (\texttt{Face Match = True}), hệ thống không chặn cứng gây ùn tắc cửa. Backend tự động chuyển nhánh Soft Mode: phát lệnh mở chốt qua MQTT, gửi payload đồng bộ đẩy dòng log mang các mã định danh chuyên biệt gồm \texttt{MISSING\_IN\_AUTO\_OPEN} hoặc \texttt{ANTI\_PASSBACK\_IN} hiển thị dạng Badge màu cam trực quan (Visual Alert) trên giao diện Web cổng 3002, đồng thời điều khiển ứng dụng trạm Android hiển thị banner thông báo đồ họa nhắc nhở nhân viên vi phạm quy trình tại chỗ.

\subsection{Kết quả thực nghiệm định lượng đạt được}

Việc ứng dụng giải pháp đa giao thức bất đồng bộ kết hợp Soft Mode mang lại hiệu quả vượt trội về mặt tối ưu hóa băng thông mạng và triệt tiêu độ trễ trạm đầu cuối. Chu kỳ thời gian xử lý trung bình cho một lượt điểm danh hai lớp khép kín ghi nhận thực tế đạt mức 3.3 giây (trong đó tiến trình bắt tay bảo mật NFC đọc chip CCCD tiêu tốn trung bình 1.5 giây, chụp ảnh selfie 0.5 giây, truyền tải payload qua môi trường internet 0.5 giây, và thuật toán trích xuất đối sánh vector đặc trưng tại máy chủ AI GPU tiêu tốn trung bình 0.8 giây). Tốc độ phản hồi hoàn toàn đáp ứng xuất sắc yêu cầu phi chức năng dưới 5 giây đặt ra tại Chương 2. Thời gian đồng bộ hiển thị màn hình Dashboard qua kênh WebSocket STOMP đạt mức lý tưởng (dưới 1 giây), chứng minh tính khả thi của mô hình Async Pipeline đề xuất.

\section{Giải pháp nâng cao năng lực an ninh hệ thống theo tiêu chuẩn OWASP Top 10}
\label{section:5.3}

\subsection{Dẫn dắt bài toán và thách thức kỹ thuật}

Một hệ thống an ninh kiểm soát ra vào sử dụng công nghệ IoT và kiến trúc phân tán luôn đối mặt với những nguy cơ tấn công mạng nghiêm trọng. Các đối tượng tấn công có thể thực hiện hành vi bắt gói tin và giả mạo chuỗi Token xác thực nhằm chiếm quyền điều khiển, cố tình gọi các endpoint nhạy cảm để leo thang đặc quyền, hoặc chèn các đoạn mã độc thông qua các trường nhập liệu tìm kiếm nhằm phá hủy toàn vẹn hệ thống cơ sở dữ liệu MySQL trung tâm. Do đó, việc thiết lập một hành lang bảo mật vững chắc cho toàn bộ hạ tầng REST API là nhiệm vụ sống còn của đồ án tốt nghiệp.

\subsection{Giải pháp kỹ thuật đề xuất}

Hệ thống tiến hành đóng gói và triển khai bộ giải pháp an ninh đa lớp bám sát mô hình tiêu chuẩn bảo mật quốc tế OWASP Top 10, thực thi đồng bộ tại lõi cấu hình của Spring Security 6:
\begin{itemize}
    \item \textbf{Phòng chống Broken Access Control qua mã hóa Token JWT phi trạng thái:} Hệ thống chèn bộ lọc tùy biến \texttt{JwtFilter} (kế thừa lớp \texttt{OncePerRequestFilter}) đứng trước bộ lọc xác thực mặc định. Mọi request gửi lên bắt buộc phải đính kèm header \texttt{Authorization: Bearer}. Chuỗi token được ký số bảo mật bằng thuật toán băm mã hóa \texttt{HS256}. Bất kỳ hành vi chỉnh sửa thủ công nào vào phần Payload của token sẽ làm sai lệch chữ ký số, máy chủ lập tức từ chối yêu cầu và trả mã lỗi HTTP 401 Unauthorized.
    \item \textbf{Phòng chống SQL Injection qua kiến trúc parameterized JPA:} Toàn bộ lớp tầng dữ liệu repository của dự án tuyệt đối không sử dụng cơ chế cộng chuỗi SQL thô. Hệ thống ứng dụng công nghệ Spring Data JPA thực hiện tham số hóa tự động (\textit{Parameterized Queries}) thông qua cơ chế PreparedStatement, bảo đảm mọi chuỗi dữ liệu độc hại nhập từ ô tìm kiếm nhân viên đều bị thoát ký tự hoàn toàn, vô hiệu hóa hoàn toàn nguy cơ tấn công chiếm quyền dữ liệu.
    \item \textbf{Kiểm soát chính sách CORS và đồng bộ phân quyền đối xứng:} Hệ thống cấu hình tường minh chính sách Cross-Origin Resource Sharing, chỉ chấp nhận đặc quyền gọi tài nguyên từ origin được chỉ định cụ thể, chặn đứng hoàn toàn các yêu cầu giả mạo từ bên ngoài. Song song với đó, tại tầng quản trị phân quyền, cơ chế Auto-create Permission đối xứng tự động tạo lập một cặp quyền truy cập (Vào/Ra) cho cùng một nhân sự tại cùng một vị trí dựa theo trường \texttt{location\_name} của bảng \texttt{devices}, đảm bảo logic State Machine đóng kín chu trình dữ liệu, triệt tiêu lỗ hổng phân quyền bất đối xứng.
\end{itemize}

\subsection{Kết quả thực nghiệm định lượng đạt được}

Hệ thống đã triển khai thực thi tổng cộng 45 kịch bản trường hợp kiểm thử hộp đen và kiểm thử tích hợp API bám sát các nguy cơ lỗ hổng an ninh thuộc danh mục OWASP. Kết quả thực nghiệm ghi nhận 43 kịch bản đạt tiêu chuẩn (đạt tỉ lệ thành công 95.56\%), các hành vi sửa đổi claims trong payload của JWT Token đều bị hệ thống bóc tách sai chữ ký và trả lỗi HTTP 401, các yêu cầu truy cập sai đặc quyền (leo thang quyền) từ tài khoản VIEWER cố gọi endpoint hành chính đều bị bộ lọc \texttt{@PreAuthorize} đánh chặn sớm và trả mã lỗi HTTP 403 Forbidden. Toàn bộ câu lệnh tìm kiếm hồ sơ nhân sự dính chuỗi SQL Injection thô đều bị tầng JPA thoát ký tự và trả về kết quả rỗng an toàn, bảo chứng cho tính sẵn sàng cao và năng lực vận hành tin cậy của sản phẩm phần mềm khi đưa vào triển khai thực tế tại các doanh nghiệp.

\clearpage

Thảo luận tổng kết chương, Chương 5 đã phân tích chuyên sâu ba giải pháp kỹ thuật nổi bật và đóng góp cốt lõi của đồ án tốt nghiệp bao gồm: kiến trúc xác thực hai lớp tự trị phân tán không kho ảnh tập trung dựa trên hạ tầng định danh chip CCCD Việt Nam, đường ống truyền thông bất đồng bộ đa giao thức tối ưu hóa băng thông mạng MQTT và WebSocket hỗ trợ cơ chế Soft Mode với chỉ báo cảnh báo trực quan (Visual Alert) trên giao diện Web cổng 3002, và hệ thống tường lửa bảo mật API đạt tiêu chuẩn an ninh OWASP Top 10. Hệ thống các thông số đo lường định lượng về độ chính xác mô hình sinh trắc học ArcFace (FAR 0.01\%, FRR 1.15\%), độ trễ luồng đầu cuối tích hợp đạt mức ấn tượng 3.3 giây đã chứng minh thành công tính khả thi vượt trội của đề tài. Từ những kết quả thực nghiệm xuất sắc và các đúc kết công nghệ toàn diện này, toàn bộ quá trình nghiên cứu, phát triển đồ án tốt nghiệp sẽ được tổng kết, đối chiếu mục tiêu và vạch ra các hướng phát triển tương lai trong chương cuối cùng — Chương 6: Kết luận và hướng phát triển.

\end{document}